%%%%%%%%%%%%%%%%%%%%%%%%%%%%%%%%%%%%%%%%%
% Medium Length Professional CV
% LaTeX Template
% Version 2.0 (8/5/13)
%
% This template has been downloaded from:
% http://www.LaTeXTemplates.com
%
% Original author:
% Rishi Shah 
%
% Important note:
% This template requires the resume.cls file to be in the same directory as the
% .tex file. The resume.cls file provides the resume style used for structuring the
% document.
%
%%%%%%%%%%%%%%%%%%%%%%%%%%%%%%%%%%%%%%%%%

%----------------------------------------------------------------------------------------
%	PACKAGES AND OTHER DOCUMENT CONFIGURATIONS
%----------------------------------------------------------------------------------------

\documentclass{resume} % Use the custom resume.cls style

\usepackage[left=0.75in,top=0.6in,right=0.75in,bottom=0.6in]{geometry} % Document margins
\usepackage{hyperref}
\usepackage{pdfpages}
\usepackage{ragged2e}
\newcommand{\tab}[1]{\hspace{.2667\textwidth}\rlap{#1}}
\newcommand{\itab}[1]{\hspace{0em}\rlap{#1}}
\name{Summer Haag} % Your name
%\address{3182 Westwood Ct, Boulder CO, 80304} % Your address
%\address{123 Pleasant Lane \\ City, State 12345} % Your secondary addess (optional)
\address{summer.haag@colorado.edu} % Your phone number and email

\begin{document}

\textbf{Research Interests:} Circle packings, Binary Quadratic Forms, arithmetic groups, continued fractions, illustrating mathematics

%----------------------------------------------------------------------------------------
%	EDUCATION SECTION
%----------------------------------------------------------------------------------------

\begin{rSection}{Education}
{\bf University of Colorado at Boulder, Boulder CO} \hfill {\em August 2022 --- Present} \\
PhD candidate \\
MA in mathematics, Fall 2024

{\bf Illustration as a Mathematical Research Technique Trimester, Paris} \hfill {\em Spring 2026}
    \begin{itemize}
        \item Three month long Trimester Program focusing on illustrating mathematics through computers, fabric, paper, etc. at Insitut Henri Poincaré along with multiple workshops
    \end{itemize}

{\bf University of Georgia, Athens GA} \hfill {\em August 2018 --- May 2022} 
\\ BA in mathematics \\
%Graduated \textit{cum laude} \\
%\\{\bf Budapest Semesters in Mathematics, Budapest Hungary} \hfill {\em August 2015 --- December 2015} 
%\\  Completed courses: \\ Number Theory, Probability, Theory of Computing, and Set Theory \\
%Minor in Linguistics \smallskip \\
%Member of Eta Kappa Nu \\
%Member of Upsilon Pi Epsilon \\


\end{rSection}

\begin{rSection}{Publications \& Preprints}

\textbf{Repellent properties of perfect powers on partition functions: a heuristic approach} \\
    S. Haag, P. Samanta, Swati, H Swisher, S Treneer, R Visser \begin{FlushRight}\textit{Research in Mathematical Science (accepted)} (\href{https://arxiv.org/abs/2601.18138}{arxiv}) \end{FlushRight}
\textbf{The local-global conjecture for Apollonian circle packings is false} \\
 	 S. Haag, C. Kertzer, J. Rickards, K. Stange \begin{FlushRight} \textit{Annals of Mathematics} (\href{https://arxiv.org/abs/2307.02749}{arxiv}, \href{https://annals.math.princeton.edu/2024/200-2/p06}{eprint}) \end{FlushRight}
\textbf{Computational study of non-unitary partitions}  \\
     A.P. Akande, T. Genao, S. Haag, M. Hendon, N. Pulagam, R. Schneider, A. Sills  \begin{FlushRight}
      \textit{J. of the Ramanujan Mathematical Society} (\href{https://arxiv.org/abs/2112.03264}{arxiv})\end{FlushRight}
\textbf{The error term in counting prime pairs}  \\
     L. Chou, S. Haag, J. Huryn, A. Ledoan \begin{FlushRight} \textit{J. of Number Theory} (\href{https://arxiv.org/abs/2308.14888}{arxiv} \href{https://doi.org/10.1016/j.jnt.2025.04.009}{eprint}) \end{FlushRight} 
     

\end{rSection}

%--------------------------------------------------------------------------------
%    Conference Talks
%-----------------------------------------------------------------------------------------------

\begin{rSection}{Invited Conference Talks}
 \textbf{Local-global Conjecture for Apollonian Circle Packings } \hfill UConn \\
 	 CT Number Theory \hfill June 2024 \\


% \textbf{Counting cycles along the spine} \hfill Howard University \\
% 	 Post-Quantum Cryptography Special Session \hfill April 2024 \\
 %	 AMS Eastern Sectional Meeting \\
\end{rSection}


%--------------------------------------------------------------------------------
%    Awards
%-----------------------------------------------------------------------------------------------
\begin{rSection}{Awards \& Nominations}


{\bf Quanta Magazine- Research Feature} \\
Featured in an article in Quanta Magazine by Max G. Levy covering the Local-Global Conjecture for
Apollonian circle packings. (\href{https://www.quantamagazine.org/two-students-unravel-a-widely-believed-math-conjecture-20230810/}{ Article reference})  \\
\newpage
{\bf Laureate of the QRT Welcome Package}\\
Received a grant of 3000€, funded by \href{https://www.qube-rt.com/}{Qube Research and Technologies} and administered by the \href{https://www.fonds-ihp.org/}{IHP Endowment fund}, set aside to promote the participation of women in IHP's scientific programs.\\
\\{\bf 2025 Throne Summer Fellowship}\\
Received 7,000 in summer funding for research\\
\\{\bf UGA Strahan Award} \\
Nominated by professors for excellence in mathematics as a junior in undergraduate.\\
\\{\bf CURO Assistantship}\\
Received 1,000 at UGA for research in Number Theory to present at the CURO Symposium.\\
\\{\bf The Hollingsworth Award}
\\ Nominated and recieved at UGA for excellence in Math.\\

\end{rSection}

%--------------------------------------------------------------------------------
%    Seminars etc
%-----------------------------------------------------------------------------------------------

\begin{rSection}{Seminars and Other Talks}

 	  \textbf{Overview of Local-to-Global Conjecture for Apollonian Circle Packings} \hfill CU Boulder\\
 	 Math for All \hfill April 2024 \\ \\ 
     \textbf{Local-to-Global Conjecture for Apollonian Circle Packings} \hfill University of Houston\\
 	 Dynamics Seminar \hfill May 2025 \\ \\
     \textbf{Local-to-Global Conjecture for Apollonian Circle Packings} \hfill Oregon State University\\
 	 Algebra and Number Theory Seminar \hfill October 2025 \\ \\
     \textbf{The Parameter Space of Apollonian Circle Packings} \hfill Aix Marseille University\\
 	 Geometry Seminar \hfill March 2026 \\ \\
\end{rSection}


%---------------------------------------------------------------------------
% Conferences and workshops - Invited
%---------------------------------------------------------------------------


\begin{rSection}{Invitational conferences and workshops}
\begin{itemize}
\item CTNT Summer School and Conference 2024 
\item Joint Mathematics Meeting 2025 --- Organizer for Special Session on Local-to-Global in Apollonian Circle Packings and Beyond. 
\item Rethinking Number Theory 6
\item Higher Dimensional Hyperbolic Geometry Workshop 2025
% \item Illustration as a Mathematical Research Technique Trimester IHP
%     \begin{itemize}
%         \item Three month long Trimester Program focusing on illustrating mathematics through computers, fabric, paper, etc. at Insitut Henri Poincaré along with multiple workshops
%     \end{itemize}
\item Illustration as a Mathematical Research Technique CIRM
\item Rigorous Illustrations - Their creation and evaluation for mathematical research
\item Bridging visualisation and understanding in Geometry and Topology
\item Integrating Research and Illustration in Number Theory
\end{itemize}

\end{rSection}




%--------------------------------------------------------------------------------
%    Conferences - addtended
%-----------------------------------------------------------------------------------------------

\begin{rSection}{Conferences attended}
\small
\begin{itemize}
\item Math for All 2023, 2024 \\
\item Integers Conference 2025 \\
\end{itemize}

\end{rSection}

%--------------------------------------------------------------------------------
%    Teaching
%-----------------------------------------------------------------------------------------------

\begin{rSection}{Teaching Experience}
{\bf Primary Instructor}
\\University of Colorado Boulder
\small
\begin{itemize}
\item Precalculus \hfill Fall 2025
\item Calculus II \hfill Fall 2023, Summer 2024, Fall 2024

\end{itemize}

{\bf Calculus Teaching Assistant}\\
University of Colorado Boulder
\small
\begin{itemize}
\item Precalculus \hfill Spring 2025
\item Calculus I \hfill Fall 2022
\item Calculus II \hfill Spring 2023
\item Calculus III \hfill Spring 2024
\end{itemize}

{\bf SageMath Workshop}\\
\small 
Introduction into using SageMath at the Integrating Research and Illustration in Number Theory workshop with worksheet written by Katerine E. Stange


\end{rSection}

%---------------------------------------------------------------------------
% Service
%---------------------------------------------------------------------------


\begin{rSection}{Service \& Mentorship}
{\bf No-Senior Seminar} \hfill \textit{Organizer} \\
A seminar at the Illustrations as a Mathematical Research Technique Trimester\\ for graduate students and Post-docs to present topics  \hfill Spring 2026. \\

{\bf Math Peer Mentorship} \hfill \textit{Mentor} \\
Older graduate students pair with first year graduate students  \hfill Fall 2023 --- Spring 2024. \\

{\bf Reading Groups}  \hfill \textit{Organizer}\\
University of Colorado, Boulder \\
\begin{itemize}
\item $p$-adic Analysis --- Summer 2024.
\end{itemize} 

{\bf Prelim Mentor Program-- Algebra} \hfill \textit{Mentor} \\
Graduate student who has taken Algebra I and II 
and passed the Algebra \\Preliminary Exam guide graduate students in studying for the \\Algebra Preliminary Exam \hfill Summer 2025 \\

{\bf Rabbit Rescue Grooming co-lead} \hfill \textit{Volunteer} \\
Rocky Mountain House Rabbit Rescue\\
Organize and train a team to groom approximately 50 rabbits \hfill Jan 2024-present \\

\small
\end{rSection}


%----------------------------------------------------------------------------------------
%	WORK EXPERIENCE SECTION
%----------------------------------------------------------------------------------------


%----------------------------------------------------------------------------------------
%	LANGUAGES AND SKILLS
%----------------------------------------------------------------------------------------

\begin{rSection}{Open Source Software}

{\bf Github} \\
Public repository of code for research in Apollonian Circle Packings \hfill \href{https://github.com/summerhaag}{https://github.com/summerhaag} \\

\end{rSection}
\end{document}